\note{3}{Thu 23 May 2024 22:34}{Fredholm Theory}
\tags{Compact operator, Fredholm operator}

\section{Fredholm Theory and Spectral Theory}


\subsection{Compact Operators}



Recall those facts:
\begin{itemize}
	\item In a complete metric space, relatively compact $\iff $ totally bounded $\iff $ sequentially compact
	\item Continuous functions preserve compact sets.
	\item (Riesz) A normed space is finite dimensional iff the unit ball is compact.
\end{itemize}

\begin{definition}[Compact Operator]
	\label{defn:CompactOperator} 
	In $\normalfont{Ban} $,
	an operator $T: E \to F$ is \textit{compact} if the following equivalent conditions hold:
	\begin{enumerate}
		\item  $T(U_{E}(0,1))$ is relatively compact, i.e. its closure is compact.
		\item  $T(B_{E}(0,1))$ is relatively compact, i.e. its closure is compact.
		\item $T$ maps bounded sets to relatively compact sets.
	\end{enumerate}
\end{definition}

\begin{intuition}
	Compact operators are ``almost finite dimensional'', just like linear operators in finite dimensional vector spaces. As we will see later, those operators form a closed two-sided ideal, so it makes sense to mod them out and form the \logseq{Calkin Algebra}.
\end{intuition}

\begin{proposition}
	\label{prop:CompactTwoSidedIdeal}
	$\mathcal{K}(X)$ is a closed two-sided ideal of $B(X)$.
\end{proposition}
\begin{proof}
	It is easy to see that $\mathcal{K}(X)$ is a linear subspace.

	Take $K \in \mathcal{K}(X)$, $T \in B(X)$. $T$ is continuous, thus preserves compact sets. Thus $TK \in \mathcal{K}(X)$. $T$ is bounded, so $KT \in \mathcal{K}(X)$.

	Finally we verify $\mathcal{K}(X)$ is a closed subspace. Let $\left( T_n \right) $ be a sequence in $\mathcal{K}(X, Y)$ and $T \in B(X, Y)$, such that $\|T_n - T\| \to 0$. We want to show $T \in \mathcal{K}(X, Y)$. For this we show $T(U(0,1))$ is totally bounded. Fix $\varepsilon > 0$, and we have $\|T_n - T\| < \varepsilon$ for large $n$. Since $T_n \in \mathcal{K}(X, Y)$, $T_n(U(0,1))$ is totally bounded, so there exists finitely many $x_1, \ldots, x_m \in U(0,1)$ such that $T_n(U(0,1)) \subset \cup_{i = 1}^m U(T_n x_i, \varepsilon)$. 

	Now take $x \in U(0,1)$, and let $x_i$ be such that $T_n x \in U(T_n x_i, \varepsilon)$. Now
	\[
		\|T x - T x_i\| 
		\le  \|T x - T_n x\| + \|T_n x - T_n x_i\| + \|T_n x_i - T x_i\|
		\le \varepsilon + \varepsilon + \varepsilon
	.\] 
\end{proof}

\begin{example}
	\begin{itemize}
		\item Every finite-rank operator is compact, since a closed bounded set in a finite-dimensional linear space is compact.
		\item Every projection into an infinite dimensional subspace is not compact, by invoking Riesz's Lemma of perpendicularity.
		\item Topological isomorphisms between infinite dimensional Banach spaces are never compact.
	\end{itemize}
\end{example}

Thus, any operator that can be approximated by finite-rank operators is compact, i.e. $\overline{\mathcal{F}(X,Y)} \subset \mathcal{K}(X,Y) $.

In a Hilbert space, the converse is also true:
\begin{proposition}
	\label{prop:CompactOperatorLimit}
	In a Hilbert space $H$, compact operators are limits of finite-rank operators. That is, $\mathcal{K}(H) = \overline{\mathcal{F}(H)} $.
\end{proposition}
\begin{proof}
	We may assume $\operatorname{dim} H = \infty $. Fix an orthonormal basis $(e_n)$. Define $P_n = \left\langle \cdot , e_1 \right\rangle e_1 + \ldots + \left\langle \cdot , e_n \right\rangle e_n$, the projection onto $\text{span}(e_1, \ldots, e_n)$. We claim that $\|T - P_n T\| \to 0$ for all $T \in \mathcal{K}(H)$.

	Fix $\varepsilon>0$. Using compactness of $T$, there exists finitely many $x_1, \ldots, x_m \in U(0,1)$ such that $T(U(0,1)) \subset \cup_{i = 1}^m U(T(x_i), \varepsilon)$. Now pick $x \in U(0,1)$ and $x_i$ be such that $T x \in U(T x_i, \varepsilon)$. We have
	\[
		\|T x - P_n T x\| \le \|T x - T x_i\| + \|T x_i - P_n T x_i\| + \|P_n T x_i - T x\|
	.\] 
	Since there are finitely many $x_i$, for large $n$ independent of $x$, we have $\|T x_i - P_n T x_i\| < \varepsilon$. Thus $\|T x - P_n T x\| < 3 \varepsilon$.
\end{proof}

This is to say that every Hilbert space has the \logseq{approximation property}. A Banach space $F$ has the approximation property if for every Banach space $E$ every compact operator from $E$ to $F$ can be approximated in the operator norm by finite-dimensional operators.

Here are some facts that we will not prove here, consult, for example, \cite{khelemskiiLecturesExercisesFunctional2006}, section 3.3.

\begin{itemize}
	\item If a Banach space has a \logseq{Schauder basis}, then it has the approximation property. In particular, the spaces $L_p(X, \mu), C(\Omega), c_0, l_p$ all have approximation property.
	\item (P. Enflo, 1972). There exist separable Banach spaces, in particular, closed subspaces in $c_0$, that do not have the approximation property. As a corollary (see Theorem 2) there are separable Banach spaces without Schauder bases.
\end{itemize}


\begin{proposition}[Schauder]
	\label{prop:SchauderCompactAdjoint}
	$T: E \to F$ is compact iff the Banach adjoint operator $T^*:F^* \to E^*$ is compact.
\end{proposition}
Warning: not to be confused with the Hilbert adjoint.
\begin{proof}
	We only need to show $T$ compact implies $T^*$ compact. The other side follows from the canonical embedding of a space into its double dual.

	Using compactness of $T$, pick $N_F \subset F$ an $\varepsilon$-net of $T U_E$. Define the functional
	\begin{align*}
		S: F^* &\longrightarrow \mathbb{C}^{N_F} \\
		f &\longmapsto S(f) = f(y_i)_{y_i \in N_F}
	.\end{align*}
	$S$ is a compact operator. Now pick $N_{F^*} \subset F^*$ such that $S N_{F^*}$ is an $\varepsilon$-net of $S U_{F^*}$. 

	Now pick $f \in U_{F^*}$ and $x \in E, \|x\| = 1$. We have for all $f \in U_{F^*}$,
	\[
		(T^* f) x = f T x \sim f y_i = \pi_i S f \sim \pi_i S f_j = f_j y_i \sim f_j T x = (T^* f_j) x
	.\] 
	The first $\sim $ is taking $y_i \in N_{F}$ such that $\|T x - y_i\| < \varepsilon$. The second $\sim $ is justified by taking $f_j \in N_{F^*}$ such that $\|S f - S f_j\| < \varepsilon$. Since $\|f\| < 1$ and $\|x\| = 1$ we have
	\[
		\|T^* f - T^* f_j\| < 3 \varepsilon
	.\] 
	That is, $T^* N_{F^*}$ is a $3 \varepsilon$-net of $T^* U_{F^*}$.
	
\end{proof}


\begin{proposition}
	\label{prop:CompactWeakTop}
	If two operators are weak topology equivalent, then they are either both compact or both not compact.
\end{proposition}

\subsubsection{Examples of Compact Operators}

\begin{proposition}
	\label{prop:DiagonalOperatorCompact}
	The diagonal operator $T_\lambda : l_p \to l_p$, $\lambda \in l_\infty $ is compact if and only if $\lambda \in c_0$ ($\lambda \to 0$ ).
\end{proposition}

\begin{proposition}
	\label{prop:ShiftOperatorNotCompact}
	The shift operators on $l_p, l_p(\mathbb{Z})$ or $L_p(\mathbb{R})$ are never compact.
\end{proposition}

\begin{proposition}
	\label{prop:MultOperatorNotCompact}
	The multiplication operator $T_f: L_p[a, b] \to L_p[a, b]$, $1 \le p \le \infty $, $f \in L_\infty [a, b]$ is not compact, unless $f = 0$.
\end{proposition}

\begin{proposition}
	\label{prop:IndefIntCompact}
	The operator of indefinite integration acting on $L_2[0,1]$, $C[0,1]$ or $L_1[0,1]$ is compact.
\end{proposition}

\begin{theorem}
	\label{prop:IntegralOperatorCompact}
	Every integral operator on $L_2[a,b]$ is compact.
\end{theorem}

\begin{proposition}
	\label{prop:DifferentiationNotCompact}
	Differentiation operators $D^k : C^{n + k}[a,b] \to C^n[a,b]$ are never compact. 
\end{proposition}

\begin{proposition}[Hilbert Adjoint Examples]
	\label{prop:ExAdjoint}
	\begin{enumerate}
		\item The adjoint operator for the diagonal operator $T_\lambda: l_2 \to l_2$ is $T_{\overline{\lambda} }$, where $\overline{\lambda} $ is complex conjugate of the sequence $\lambda$.
		\item The adjoint operator for the left shift operator $T_l: l_2 \to l_2$ is the right shift operator $T_r: l_2 \to l_2$, and vice versa.
		\item The adjoint operator for the multiplication operator $T_f: L_2(X, \mu) \to L_2(X, \mu)$ is $T_{\overline{f} }$.
		\item The adjoint operator for the indefinite integration operator on $L_2[0,1]$, $T: x \mapsto \int 0^t x(s) ds$, is $T^*: x \mapsto \int _t^1 x(s) d s$.
		\item Let $T_K:L_2[a,b] \to L_2[a,b]$ be an integral operator with the kernel $K(s, t)$. Then the adjoint operator is the integral operator $T_{\overline{K} }$.
	\end{enumerate}
\end{proposition}

\subsection{Fredholm Operators}

On a Banach space, since $\mathcal{K}(X,Y)$ is closed two-sided ideal of $B(X,Y)$, we can form the category $\mathbf{Ban}/\mathcal{K}$, where the objects are same as in $\mathbf{Ban}$, and the morphisms $\mathbf{h}_{\mathbf{Ban}/\mathcal{K}}$ are cosets $f + \mathcal{K}(E,F)$ where $f$ is in $\mathbf{h}_{\mathbf{Ban}}$.

In $\mathbf{Ban}/\mathcal{K}$, the zero maps are exactly the compact operators. The isomorphisms are, in some sense, ``anti-compact'' operators, and we are going to see that they are precisely cosets of \textit{Fredholm operators}.

This idea is more straightforward in the case of bounded operators on Hilbert spaces, and we are going to develop some theory about it later on.

\begin{definition}[Calkin Algebra]
	\label{defn:CalkinAlgebra}
	
	On a Hilbert space $H$, the \textit{Calkin algebra} is the quotient $C^*$-algebra $B(H) / K(H)$, satisfying the following short exact sequence in the category of $C^*$-algebras:
	\[
		0 \to K(H) \to B(H) \to B(H)/K(H) \to 0
	.\] 
\end{definition}

Back to Banach spaces, let's first make the central definition.

\begin{definition}[Fredholm Operator]
	\label{defn:FredholmOperator}
	In $\mathbf{Ban}$, $S: E \to F$ is a \textit{Fredholm operator} if  $\operatorname{dim} \operatorname{ker} S < \infty $ and $\operatorname{dim} \operatorname{coker} S < \infty $. 

	The \textit{index} of $S$ is $\operatorname{ind } S = \operatorname{dim} \operatorname{ker} S - \operatorname{dim} \operatorname{coker} S$.
\end{definition}

The most important facts we are going to justify are:
\begin{itemize}
	%\item $T \in \mathcal{K}(E,F)$ is never Fredholm.
	%\item For $\lambda \neq 0 \in \mathbb{C}$, $\lambda1 - T$ is Fredholm.
	\item For $S: E \to F$ and $R: F \to G$ Fredholm, $\operatorname{ind}(RS) = \operatorname{ind } R + \operatorname{ind } S $.
	\item $S$ is Fredholm iff its coset $S + \mathcal{K}(E, F)$ is an isomorphism in $\mathbf{Ban}/\mathcal{K}$.
	\item $\Phi(E, F)$ is open in $B(E, F)$ and  $S \mapsto \operatorname{ind } S$ is continuous.
\end{itemize}

\begin{example}
	\begin{itemize}
		\item In finite-dimensional vector spaces, for a linear map $T: E \to F$, we always have $\operatorname{dim} \operatorname{ker} T + \operatorname{dim} \operatorname{Im} T  = \operatorname{dim} E$, which implies $\operatorname{Ind } T = \operatorname{dim} \operatorname{ker} T - (\operatorname{dim} F - \operatorname{dim} \operatorname{Im} T) = \operatorname{dim} E - \operatorname{dim} F$. In infinite dimensions, however, the index of $T$ will not follow simply from the geometry of the spaces, because $\operatorname{Im} T$ may be infinite.
		\item Every topological isomorphism between Banach spaces is Fredholm with index $0$.
		\item Every bounded projection with finite codimension $n$ is Fredholm with index $-n$.
		\item The left and right shift operators in $l_p(\mathbb{N})$ are Fredholm with index $1$ and $-1$.
		\item The differential operator $T: x\mapsto x^{(n)} + f_1(t) x^{(n - 1)} + \ldots + f_{n - 1} x^{(1)} + f_n(t)$, $f_1, \ldots, f_n \in C[a,b]$ from $C^n[a,b]$ to $C[a,b]$ has index $n$.
		\item The derivative $x \mapsto x'$ from $C^1(\mathbb{T})$ to $C(\mathbb{T})$ has index $0$.
	\end{itemize}
\end{example}
\begin{proof}(Of last two statements.)
	$D^n: C^n[a,b] \to C[a,b]$ has index $n$ since $\operatorname{ker} D^n = \{f: D^{n - 1} f \text{ is constant}\} $, we may do indefinite integration to $D^n f = 0$ to get $f = c_1 + c_2 x + \ldots + c_n x^{n - 1}$.

	Now verify for $T$. The general theory of ODE says that there are $n$ independent solutions to $T x = 0$, i.e. $\operatorname{dim} \operatorname{ker} T = n$. The cokernel is 0 since $\operatorname{Im} T$ is dense in $C[a,b]$; $x^{(1)}, \ldots, x^{(n)}$ can be independently tweaked to cover the whole range of $T$.

	Verify $D: C^1[\mathbb{T}] \to C[\mathbb{T}]$ has index 0: we have $\operatorname{ker} D = \{f: \text{constant on } \mathbb{T} \} = \mathbb{C} $, and $\operatorname{Im} D = \{f: \int f = 0\} $ by fundamental theorem of calculus. Then the sequence
	\[
		C^1[\mathbb{T}] \xrightarrow{D} C[\mathbb{T}] \xrightarrow{\int } \mathbb{C} 
	\] 
	is exact. From this we know $\operatorname{coker} D = \mathbb{C}$, so $\operatorname{ker} D = \operatorname{coker} D$.
\end{proof}


\begin{nonexample}
	\begin{proposition}
		\label{prop:CompactNotFred}
	In the case of infinite dimensional $E$ or $F$, compact operators $T: E \to F$ are never Fredholm.
	\end{proposition}

	\begin{proof}
		If $T$ is Fredholm, then $\operatorname{Im} T$ is a Banach space. Suppose $F$ infinite dimensional, then $\tilde T: E \to \operatorname{Im} T$ is a surjective operator onto an infinite-dimensional Banach space, hence $T(U(0,1))$ is open in $\operatorname{Im} T$, which cannot be compact (since unit ball in infinite dimensional normed space is not compact).

		If $F$ is finite dimensional, then  $\operatorname{ker} T$ finite dimensional implies $E$ finite dimensional, a contradiction.
	\end{proof}
\end{nonexample}

The index is a better property than the $\operatorname{dim} \operatorname{ker} T$ and $\operatorname{dim} \operatorname{coker} T$, since it satisfies the multiplicative property.

\begin{theorem}[Multiplicative property of index]
	\label{thm:MultiplicativePropertyIndex}
	Let $S: E \to F$ and $R: F \to G$ Fredholm, then $RS$ is Fredholm, and $\operatorname{ind}(RS) = \operatorname{ind } R + \operatorname{ind } S $.
\end{theorem}
\begin{proof}
	It's easily checked that $RS$ is Fredholm. To compute index, consider the following exact sequence:
	\[
		0 \to \operatorname{ker} S \hookrightarrow \operatorname{ker} RS \xrightarrow{S} \operatorname{ker} R \xrightarrow{+ \operatorname{Im} S} \operatorname{coker} S 
		\xrightarrow{R} \operatorname{coker} RS \xrightarrow{+ \operatorname{Im} R} \operatorname{coker} R \to 0
	.\] 
	The desired equation follows from the following fact:
\begin{lemma}
	\label{lem:EulerCharacteristicZero}
	The \logseq{Euler Characteristic} of an exact sequence in the category of finite dimensional vector spaces is zero. To be specific, let
\[
0 \rightarrow A^0 \xrightarrow{d_0} A^1 \xrightarrow{d_1} A^2 \xrightarrow{d_2} \cdots \rightarrow A^m \rightarrow 0
\]
be an exact sequence of finite-dimensional vector spaces. Then
\[
\sum_{k=0}^m(-1)^k \operatorname{dim} A^k=0
\]
\end{lemma}
The Lemma follows from the rank-nullity theorem from linear algebra,
\[
\operatorname{dim} A^k=\operatorname{dimker} d_k+\operatorname{dimim} d_k .
\]
Use exactness condition to simplify the sum.

\end{proof}

Next, we study the stability of index under compact perturbation.
For $T: E \to E$ compact, consider the operator $S = 1 - T$. This is an important special class of Fredholm operators, originally motivated by the study of \textit{integral equations} $f - \int K f = g$.

\begin{theorem}
	\label{thm:ITFredholm}
	$S = 1 - T$ is Fredholm.
\end{theorem}
\begin{proof}
	\textbf{Kernel is finite dimensional.} For each $x \in \operatorname{ker} S$, we have $x = T x$. Thus $T|_{\operatorname{ker} S}$ is both an identity operator, and a compact operator (since $T$ is). Hence $\operatorname{ker} S$ is finite dimensional. We can then take $E_S$ be a closed subspace that is linear complement of $\operatorname{ker} S$, i.e. $E = \operatorname{ker}  S \oplus E_S$.

	\textbf{Image is closed.} Take $x_n' \in E$ such that $y_n = S x_n' \to y \in E$. We want to show $y \in \operatorname{Im} S$. Let $x_n \in E_S$ be projection of $x_n'$.

	We first verify $(x_n)$ is bounded. Suppose in contrary that $\|x_n\| \to \infty $. Then $x_n^0 = \frac{x_n}{\|x_n\|}$ are unit vectors such that $S x_n^0 \to 0$. Since $T$ is compact, $T x_n^0 \to z$ after passing to a subsequence. Since $\|x_n^0 - T x_n^0\| \to 0$, $x_n^0 \to z$ along the same subsequence. Thus $S z = 0$. But now we have  $z \in \operatorname{ker} S, z \in E_S, $ and $\|z\| = 1$, which is impossible.

	Since $(x_n)$ is bounded, $T x_n \to z$ after passing to a subsequence. We have $y_n = S x_n = x_n - T x_n$, so $x_n = y_n + T x_n \to y + z$, whence $S x_n \to S(y + z) = y$.

	\textbf{Image has finite codimension.} First observe that $\operatorname{Im} S$ is an invariant space of $S$, and $T$ as well. (Since $T S = S T$, we have $T (S x) = S(T x)$ for all $x \in E$) 
	We can form the quotient $\tilde E = E / \operatorname{Im} S$, and have operators $\tilde S, \tilde T: \tilde E \to \tilde E$. Now for every $\dot{x} \in \tilde E$, we have $\tilde S \dot{x} = 0$ and $\tilde S \dot{x} = \dot{x} - \tilde T \dot{x}$, so $\tilde T$ is identity on $\tilde E$.

	Let $\operatorname{pr} : E \to \tilde E$ be the natural projection. Then
	\[
		(\operatorname{pr}  \circ T)(B_E) = (\tilde T \circ \operatorname{pr} )(B_E) = \tilde T (B_{\tilde E}) = B_{\tilde E}
	.\] 
	Since $\operatorname{pr }  \circ T$ is compact ($\mathcal{K}$ (E) is two-sided ideal), $B_{\tilde E}$ is relatively compact. Hence $\operatorname{dim} \tilde E < \infty $, i.e. $\operatorname{dim} \operatorname{coker} S < \infty $.
\end{proof}

\iffalse
\begin{proof}[Alternative proof for finite codimension]
Suppose for contradiction that  $\operatorname{Im} S$ had infinite codimension. Then we can pick a sequence of strictly increasing subspaces $\operatorname{Im} S  = V_0 \subsetneq V_1 \subsetneq \ldots$, increasing dimension one at a time. Using the Riesz Lemma ( Lemma~\ref{lem:Riesz}), we can choose unitary $x_n \in V_n, \text{dist}(x_n, V_{n - 1}) > \frac{1}{2}$. But $S x_n = x_n - T x_n \in V_{n - 1}$. 
\end{proof}
\fi

\begin{theorem}[Nikol'skii, Atkinson]
	\label{thm:NikolskiiAtkinson}
	Consider $S : E \to F$ in $\mathbf{Ban}$. Then $S$ is a Fredholm operator $\iff $ $S + \mathcal{K}(E, F)$ is an isomorphism in $\mathbf{Ban}/\mathcal{K}$. In other words, $S$ is Fredholm $\iff $ there exists $R: F \to E$ such that $RS = 1_E - T_1, SR = 1_F - T_2$ where $T_1 \in \mathcal{K}(E), T_2 \in \mathcal{K}(F)$.

	Moreover, $R$ can be chosen such that $T_1$ and $T_2$ are finite-dimensional.
	
\end{theorem}
\begin{proof}
	($\implies$)
	Take Fredholm $S: E \to F$. Since  $\operatorname{dim} \operatorname{ker} S < \infty $, we can take its closed linear complement $E_S$, so there is a projection $P_1$ onto $E_S$ along $\operatorname{ker} S$. Since $\operatorname{dim} \operatorname{coker} S < \infty $, $\operatorname{Im} S$ is a closed subspace and there is a projection $P_2$ onto $\operatorname{Im} S$. 

	Let $S_0^0: E_S \to \operatorname{Im} S$ be the corestriction of $S$, then it is a bijective map between Banach spaces. Hence by Banach theorem (open mapping theorem) it is a topological isomorphism. 

	Now let $R: F \to E$ be defined as $R = (S_0^0)^{-1} P_2$, where $E_S$ naturally embeds into $E$. We have $RS = (S_0^0)^{-1} P_2 S_0^0 P_1 = P_1 = 1 - (1 - P_1)$ and $SR = S_0^0 P_1 (S_0^0)^{-1} P_2 = P_2 = 1 - (1 - P_2)$, where both $1 - P_1$ and $1 - P_2$ are finite dimensional (since they are projections onto $\operatorname{ker} S$ and $\operatorname{coker} S$).

	($\impliedby$) From the assumption, by Theorem~\ref{thm:ITFredholm} we know $RS$ and $SR$ are Fredholm operators. Then observe that $\operatorname{ker} S \subset \operatorname{ker} RS$ and $\operatorname{Im} S \supset \operatorname{Im} SR$.

\end{proof}

\begin{theorem}[Continuity of Index]
	\label{thm:ContinuityIndex}
	$\text{Fred}(E,F)$ is an open subset of $B(E,F)$. The index map $S \mapsto \operatorname{ind } S$ is locally constant and in particular continuous.
\end{theorem}
\begin{proof}
	By Theorem~\ref{thm:NikolskiiAtkinson}, there are closed subspaces $E_0 \subset E$ and $F_0 \subset F$, each with finite codimension, such that the corestriction $S_0^0: E_0 \to F_0$ is a topological isomorphism. Write $\iota : E_0 \to E$ for the natural embedding and $\pi : F \to F_0$ for the natural projection. Then $S_0^0 = \pi \circ S \circ \iota $. 

Now, $\text{Iso}(E_0, F_0)$ is open in $B(E_0, F_0)$, since $\text{Iso}(E_0, E_0) = \text{GL}(B(E_0))$ is open (Proposition~\ref{prop:GLOpen}), and for each $V \in \text{Iso}(F_0, E_0)$,
	 \[
		\text{Iso}(E_0, F_0) \ni T \mapsto VT \in \text{Iso}(E_0, E_0)
	\] 
	is a homeomorphism.

	Now, choose $\varepsilon_0>0$ such that $U(\pi S \iota, \varepsilon_0) \subset \text{Iso}(E_0, F_0) $. Now, fix $S \in \text{Fred}(E, F)$ and let $R$ be such that $\|R - S\| < \frac{\varepsilon_0}{(\|\iota\|+1)(\|\pi\|+1)}$, then
	\[
		\|\pi R \iota - \pi S \iota\|
		\le \|\pi\| \|R - S\| \|\iota\|
		< \varepsilon_0
	.\] 
	Thus  $\operatorname{ker} R \subset \operatorname{ker} S, \operatorname{Im} R \supset \operatorname{Im} S$, thus $R \in \text{Fred}(E, F)$. Now both $\pi R \iota$ and $\pi S \iota$ are topological isomorphisms with Fredholm index $0$, so by Theorem~\ref{thm:MultiplicativePropertyIndex} we have
	\[
		\text{ind}\pi + \text{ind}R + \text{ind}\iota =
		\text{ind}\pi + \text{ind}S + \text{ind}\iota = 0
	.\] 
	Thus $\text{ind} R = \text{ind} S$.
\end{proof}

\begin{corollary}[Stability of index under compact perturbations]
	\label{cor:CompactPerturbation}
	If $S$ is a Fredholm operator and $T$ a compact operator between Banach spaces $E$ and $F$, then $S + T$ is a Fredholm operator and $\text{ind}(S + T) = \text{ind}(S)$.
\end{corollary}
\begin{proof}
	The map $[0,1] \ni \text{ind}(T + t K) \in \mathbb{Z}$ is continuous.
\end{proof}

\begin{corollary}
	\label{cor:ITindex}
	If $K \in \mathcal{K}(E)$, the operator $S = 1 - K$ has index $0$.
\end{corollary}





\begin{theorem}[Fredholm Alternative]
	\label{thm:FredholmAlternative}
	Let $T \in \mathcal{K}(E)$ be a compact operator. Then for each $\lambda \in \mathbb{C}$, $\lambda\neq 0$, either $\lambda$ is an eigenvalue of $T$ with finite dimensional eigenspace, or $\lambda - T$ invertible.
\end{theorem}

Other equivalent statements of Fredholm Alternative:
\begin{itemize}
	\item 
	either $0 < \operatorname{dim} \operatorname{ker} 1 - T < \infty $, or $1 - T$ invertible. 
	\item $1 - T$ injective $\iff  1 - T $ surjective.
	\item either $(\lambda - T) v = 0$ has some nonzero solution $v \in E$, or $(\lambda - T) v = w$ always have unique solution $v$ for arbitrary  $w \in E$, in which case $v$ depends continuously on $w$.
\end{itemize}

\begin{proof}
	$\lambda \mathbf{1}$ is a Fredholm operator. By the above corollary, $\lambda - T$ is Fredholm with index $0$. Then either $\operatorname{dim} \operatorname{ker} (\lambda - T) = \operatorname{dim} \operatorname{coker} (\lambda - T) 0$, in which case $\lambda - T$ is invertible, or $0 < \operatorname{ker} T < \infty $, in which case $\lambda$ is an eigenvalue.
\end{proof}

We have seen from Theorem~\ref{thm:ContinuityIndex} that $\Phi(E, F)$ is union of open-closed subsets $\Phi_n(E, F)$ of Fredholm operators with index $n$, $n \in \mathbb{Z}$.

\begin{theorem}[Atiyah]
	\label{thm:AtiyahFredConnected}
	Let $H$ and $K$ be Hilbert spaces. Then every $\Phi_n(H, K)$ is non-empty and path-connected.
\end{theorem}

\subsubsection{Homological Algebra and Adjointness}

\begin{proposition}
	\label{prop:FredholmExactSequence}
	In category $\mathbf{Ban}$, an operator $S: E \to F$ is Fredholm iff there is an exact sequence
	\[
		0 \to K \to E \xrightarrow{S} F \to C \to 0
	.\] 
	with finite-dimensional $K$ and $C$.
\end{proposition}
\begin{proof}
	Omitted due to simplicity.
\end{proof}


\begin{theorem}[Banach Adjoint is Exact]
	\label{thm:BanachAdjointExact}
	The Banach adjoint functor is exact.
\end{theorem}
\begin{proof}
	\todo{Not done yet.}
\end{proof}


\begin{proposition}
	\label{prop:FredholmAdjoint}
	In $\mathbf{Ban}$, $S: E \to F$ is Fredholm iff $S^*: F^* \to E^*$ is Fredholm. Furthermore, we have $\operatorname{dim} \operatorname{ker} S^* = \operatorname{dim} \operatorname{coker} S$, $\operatorname{dim} \operatorname{coker} S^* = \operatorname{dim} \operatorname{ker} S$, and their index are the same.
\end{proposition}
\begin{proof}
	$S$ Fredholm implies $S^*$ Fredholm by applying Theorem~\ref{thm:BanachAdjointExact} to Proposition~\ref{prop:FredholmExactSequence}, noting that for finite dimensional $C$ and $K$, $C \cong C^*$ and $K \cong K^*$. 

	Now assume $S^*$ Fredholm. If we can show $\operatorname{Im} S$ is closed, then applying the first part to $S$ again gives the result. To do this, consider the natural embedding $\alpha: E\to E^{* *}, F \to F^{* *}$. $S^*$ Fredholm implies $S^{* *}: E^{* *} \to F^{* *}$ Fredholm. We have the following commutative diagram:
% https://quiver.local:8080/#q=WzAsNCxbMCwwLCJFXnsqICp9Il0sWzIsMCwiRl57KiAqfSJdLFswLDIsIkUiXSxbMiwyLCJGIl0sWzAsMSwiU157KiAqfSJdLFsyLDMsIlMiXSxbMywxLCJcXGFscGhhIiwxLHsic3R5bGUiOnsidGFpbCI6eyJuYW1lIjoiaG9vayIsInNpZGUiOiJ0b3AifX19XSxbMiwwLCJcXGFscGhhIiwxLHsic3R5bGUiOnsidGFpbCI6eyJuYW1lIjoiaG9vayIsInNpZGUiOiJ0b3AifX19XV0=
\[\begin{tikzcd}[ampersand replacement=\&]
	{E^{* *}} \&\& {F^{* *}} \\
	\\
	E \&\& F
	\arrow["{S^{* *}}", from=1-1, to=1-3]
	\arrow["S", from=3-1, to=3-3]
	\arrow["\alpha"{description}, hook, from=3-3, to=1-3]
	\arrow["\alpha"{description}, hook, from=3-1, to=1-1]
\end{tikzcd}\]
We may omit $\alpha$ without loss of rigor.

We have $\operatorname{ker} S \subset \operatorname{ker} S^{* *}$ finite dimensional, so  we may take its orthogonal complement $E_0$ in $E$ and $E^1$ in $E^{* *}$. Thus  $E = K \oplus E_0$ and $E^{* *} = K \oplus E_1$. In particular, $\operatorname{Im} S^{* *}$ is closed.

To show  $\operatorname{Im} S$ closed, we show $\alpha \operatorname{Im} S = S^{* *} ( \alpha E) = S^{* *}(\alpha E_0)$ is closed. Take a sequence $x_n \in E_0$ such that $S^{* *} x_n \to y \in F^{* *}$. Since $S^{* *}$ is Fredholm, $S^{* *}$ after corestricted to $E_1$ and $\operatorname{Im} S^{* *}$ is a topological isomorphism.
Then the Banach theorem says $x_n \to x \in E_1$. But  since $E_0$ is closed, $x \in E_0$, and we have $S^{* *} x = y$.

%Suppose $\|x_n\| \to \infty $. Then we let $x_n^0 = \frac{x_n}{\|x_n\|}$, and $S^{* * } x_n^0 \to 0$. Since $S^{* *}$ is Fredholm, $S^{* *}$ after corestricted to $E_1$ and $\operatorname{Im} S^{* *}$ is a topological isomorphism, thus $x_n^0 \to 0$, a contradiction. Thus $(x_n)$ is bounded. 

	
\end{proof}

\begin{remark}
	If we work in the category of Hilbert spaces instead, then $T^{* *} = T$ and we only keep the first line of the proof.
\end{remark}


\subsubsection{Examples on Lp and lp spaces}

\begin{proposition}
	A diagonal operator $T_\lambda : l_p \to l_p$, $1 \le p \le \infty $ is Fredholm iff zero is not a limit point of the sequence $\lambda$. If it is Fredholm, its index is always equal to zero.
\end{proposition}
\begin{proof}
	If $T_\lambda$ is Fredholm, its birestriction to the (infinite-dimensional) subspace generated by arbitrary family of unit vectors is still Fredholm, thus such birestrictions cannot be compact. If zero is a limit point of the sequence $\lambda$, we may choose sequence of unit vectors $\delta_{n_k} \in l_p$ such that $\lambda_{n_k} < \frac{1}{k}$. Then $T_\lambda$ restricted to subspace spanned by $(\delta_{n_k})$, is compact, since the Hilbert cube is compact.

	If zero is not a limit point, then without loss of generality we may find a lower bound $\theta>0$ of all $\lambda_n$. Then $T_\lambda$ becomes invertible, hence Fredholm.

To see $\text{ind} T_\lambda = 0$, it suffices to observe that both $\operatorname{ker} T_\lambda$ and $\operatorname{coker} T_\lambda$ are the space $\operatorname{span } \{\delta_k : \lambda_k = 0\}$.
\end{proof}

\begin{proposition}
	\label{prop:LpMultFredholm}
	The multiplication vector $T_f: L_p[a,b] \to L_p[a,b]; 1 \le p \le \infty $, where $f \in L_\infty [a,b]$ is Fredholm $\iff $ $\|f\| > \theta > 0$ a.e. on $[a,b]  \iff $ $T_f$ is invertible.
\end{proposition}
\begin{proof}
	We first prove $T_f$ invertible $\iff $ $T_f$ Fredholm. We only need to prove $\impliedby$. In case $K_0 = \{|f| = 0\} $ as positive measure, $\operatorname{ker} T_f = L(K_0)$ has infinite dimension, a contradiction. Thus $K_0$ has zero measure, then $T_f$ is injective. Since $T_f$ Fredholm, $\operatorname{Im} T_f$ is closed (Lemma~\ref{lem:ImageClosed}). Thus $T_f$ is topologically injective, thus $T_f$ is invertible (see the following lemma). 
\end{proof}

\begin{lemma}
	The following are equivalent:
	\begin{enumerate}
		\item $T_f$ is a topological isomorphism.
		\item $T_f$ is topologically injective.
		\item $T_f$ is topologically surjective.
		\item $T_f$ is surjective.
		\item $|f| > \theta > 0$.
	\end{enumerate}
\end{lemma}
\begin{proof}
	Assume $5$, then $T_{f^{-1}} = T_f^{-1}$ is bounded. Thus $T_f$ is topological isomorphism.

	Assume $5$ is false. Let $K_n = \{|f| < \frac{1}{n}\} $. Write $\chi_n = \chi_{K_n}$. Consider
	\[
		T_f \frac{\chi_n}{\|\chi_n\|} = \frac{f \chi_n}{ \|\chi_n\|} < \frac{1}{n} \frac{\chi_n}{\|\chi_n\|} \to 0
	.\] 
	Thus $T_f$ is not bounded below, thus not invertible.

	$1 \to 2, 3, 4$ is trivial. Assume $2$. Topologically injective means $T_f U(0,1)$ is open, thus $T_f$ surjective, thus $T_f$ is topological isomorphism (See Lemma).

	Assume 4. Then $1$ follows by Open Mapping Theorem.
\end{proof}

\begin{lemma}
	\label{lem:TopInjectiveSurjective}
	In the category of normed spaces $\mathbf{Nor}$, for $T:E \to F$, TFAE:
	\begin{enumerate}
		\item $T$ is topological isomorphism,
		\item $T$ is topologically injective and surjective,
		\item $T$ is injective and topologically surjective
	\end{enumerate}
\end{lemma}

\subsubsection{Motivations for Fredholm Operators}

The study of compact and Fredholm operators originally stem from the study of \logseq{integral equations of the second kind}, of the form
\[
	x(s) - \int _a^b K(s, t) x(t) dt = y(s)
.\] 
When $y(s) = 0$ this is a \textit{homogeneous} integral equation. From algebraic point of view this is an \textit{operator equation}. We consider $E = F = L_2[a,b]$, and $S = 1 - T$, where $T$ is an integral operator on $L_2[a,b] $ with kernel $K(s, t)$. 

The theory of Fredholm operators gives us the following, classic result:

\begin{theorem}[Triple Fredholm Theorem]
	\label{thm:TripleFredholmThm}
	Consider the integral equation of the second kind
\[
	x(s) - \int _a^b K(s, t) x(t) dt = y(s)
.\]
Then
\begin{enumerate}
	\item There is a finite family $x_1(s), \ldots, x_m(s)$ of linearly independent and spanning solutions to the homogeneous equation.
	\item There is a finite family $z_1(s), \ldots, z_n(s)$ of linearly independent functions in $L_2[a,b]$ such that the right hand sides $y(s)$ of which the inhomogeneous equation has at least one solution are precisely those functions of which $\int _a^b y(t) \overline{z_k(t)} dt = 0 $ for $k = 1, \ldots, n$.
	\item $m = n$.
\end{enumerate}

\end{theorem}
\begin{proof}
	In the Hilbert space $H = L_2[a,b]$, consider the operator $S = 1 - T$, where $T$ is the integral operator with kernel $K(s, t)$. Part (1) is equivalent to saying  $\operatorname{dim} \operatorname{ker} S  < \infty $, a direct corollary of Theorem~\ref{prop:IntegralOperatorCompact} and Theorem~\ref{thm:ITFredholm}. We have $m = \operatorname{dim} \operatorname{ker} S$.

	(2) is equivalent to
	\[
		\operatorname{Im} S = \left( \operatorname{span } \{z_1(s), \ldots, z_n(s)\}  \right) ^\perp 
	.\] 
	This is equivalent to saying that $\operatorname{Im} S$ is closed and $\operatorname{dim} (\operatorname{Im} S) ^{\perp} < \infty $. This is true since $\operatorname{dim} \operatorname{Im} (S) ^\perp =\operatorname{dim} \operatorname{coker} S < \infty  $. Thus $n = \operatorname{dim} \operatorname{coker} S$.

	Finally, Corollary~\ref{cor:ITindex} implies $m = n$.
\end{proof}

In the self-adjoint case Theorem~\ref{thm:TripleFredholmThm} has a simplification.
\begin{proposition}[Fredholm]
	\label{prop:FredholmSA}
	Suppose
	\[
		x(s)-\int_a^b K(s, t) x(t) d t=y(s)
	\]
	is an integral equation of the second kind, with symmetric kernel, and 
\[
		x(s)-\int_a^b K(s, t) x(t) d t=0
	\]
	is the corresponding homogeneous equation. Then there exists a linearly independent system $x_1, \ldots, x_m$ of solutions of the homogeneous equation satisfying the following conditions:
	\begin{enumerate}
		\item every solution of the homogeneous equation is a linear combination of $x_1, \ldots, x_m$;
		\item the $y(s)$ for which the inhomogeneous equation has at least one solution are precisely the functions satisfying the equations $\int _a^b y(t) \overline{x_k(t)} dt = 0 $ for all $k = 1, \ldots, m$.
	\end{enumerate}
\end{proposition}
\begin{proof}
	Let $S = 1 - T_K$ where  $T_K$ is the integral operator with kernel $K$. Then $S$ is self-adjoint. We have
	$\operatorname{ker} S = \operatorname{ker} S^* = \left( \operatorname{Im} S \right) ^\perp$.
\end{proof}






